\section{Introduction}
\label{sec:introduction}

LLM-based automated program repair systems increasingly rely on test-time compute scaling: generating $K$ candidate patches and selecting the best via a verifier~\cite{snell2024scaling}. Full execution verification is reliable but expensive (\$1--3 per candidate at repository scale\footnote{Based on per-instance costs reported by \citet{yang2024sweagent} and \citet{xia2024agentless}.}), driving demand for execution-free alternatives. Three distinct formulations have emerged for this task: \textbf{generation-based simulation}~\cite{meta2025cwm,maimon2026selfexecution} predicts execution output strings per test; \textbf{trajectory-level scoring}~\cite{shum2025swerm} predicts a binary patch verdict without per-test decomposition; and \textbf{per-test classification}~\cite{ni2023lever,chen2023codet} predicts a pass/fail token per test. Each makes different tradeoffs between diagnostic granularity, output complexity, and inference cost.

Despite active development along each axis, no prior work provides a controlled comparison holding model capacity, training data, and compute constant. Practitioners choosing a formulation for their verification pipeline must rely on results from different models, different datasets, and different evaluation protocols. Whether one formulation is uniformly preferable, or whether the choice depends on properties of the code being verified, remains an open empirical question. A related question is whether formulation choice matters at all relative to more basic decisions: at what complexity level does fine-tuning become necessary, and does formulation selection add value beyond what fine-tuning alone provides?

We present the first systematic comparison of these three formulations across two model families (Qwen3, DeepSeek-Coder), two complexity levels (function-level and repository-level), and multiple seeds. Our results reveal a three-layer structure in the verification design space. Repository-level verification is inherently harder than function-level (0-shot accuracy drops below majority baseline). Fine-tuning is necessary under standard prompting (0-shot$\to$fine-tuned gains of 25--35\pp{} at repo level). Even after fine-tuning, \emph{formulation choice is governed by a formulation $\times$ model $\times$ complexity interaction}: at function-level complexity, all three formulations achieve statistically equivalent accuracy; at repository-level complexity, formulations diverge sharply, with the divergence being formulation-inherent in direction but model-contingent in magnitude. Figure~\ref{fig:reformulation} illustrates the three formulations.

We make the following contributions:

\begin{itemize}[leftmargin=*,itemsep=2pt]
\item \textbf{Formulation equivalence at low complexity.} On function-level benchmarks, \cls{}, \cwm{}, and \traj{} achieve near-identical accuracy across two model families (Qwen3-4B: 88.4/87.2/86.3\%, 5 seeds; DeepSeek-6.7B: 86.9/86.8\%, 2 seeds). The \cls{}--\cwm{} gap is 0.02--1.24\pp{}, within seed noise in both cases. This null result establishes that classification is not inherently easier than generation when execution outputs are short and structured.

\item \textbf{Complexity-dependent divergence with formulation $\times$ model interaction.} At repository-level complexity, \cls{} achieves 19--26\pp{} higher accuracy than \cwm{} on Qwen3. Cross-model validation with DeepSeek-Coder confirms the divergence is formulation-inherent: both families show \cwm{} $\gg$ \cls{} degradation from function to repo level (Qwen3: ${\sim}$5/22\pp{}; DeepSeek: ${\sim}$12/25\pp{} for \cls{}/\cwm{}). The magnitude is model-contingent: DeepSeek exhibits larger \cls{} degradation (${\sim}$12\pp{} vs.\ ${\sim}$5\pp{}). %%SCENARIO-A%% Beyond accuracy, \cwm{} exhibits higher seed sensitivity at repo level (${\sim}$5\pp{} across seeds vs.\ ${\sim}$2\pp{} for \cls{}), while showing no overtraining under our protocol. %%END-SCENARIO-A%%

\item \textbf{Error complementarity and ensemble potential.} \cls{} and \cwm{} exhibit systematically opposite error patterns at function level: \cls{} over-predicts passing (75.5\% of errors are false negatives) while \cwm{} over-predicts failure (83.1\% false positives). An oracle ensemble achieves 90.63\% (+3.12\pp{} over the best single formulation). The formulations appear to capture complementary aspects of code behavior.

\item \textbf{Diagnostic granularity at equivalent accuracy.} Per-test classification achieves statistically equivalent variant-level accuracy to trajectory-level scoring (TOST equivalence $p = 0.007$, $\pm$2\pp{} margin) while providing per-test diagnostic information at 217$\times$ higher throughput than generation (1 output token vs.\ avg 477). Whether this additional granularity translates to downstream improvements (e.g., targeted re-repair) remains an open question.
\end{itemize}
